\documentclass[10pt,a4paper]{amsart}
\usepackage[margin=19mm]{geometry}
\usepackage{amsmath,amssymb,mathtools}
\usepackage[scr=rsfs,cal=euler]{mathalfa}
\usepackage{xfrac}
\usepackage[unicode,hidelinks,bookmarks=false]{hyperref}
\newcommand{\ie}{i.e.\ }
\newcommand{\cf}{cf.\ }
\newcommand{\resp}{respectively\ }
\newcommand{\Subsection}{\subsection}
\providecommand{\nobreakdash}{\nobreak\hbox{-}}
\setlength{\emergencystretch}{2em}
\multlinegap0pt
\usepackage{bm}
\usepackage{bbm}
\usepackage{dsfont}
\usepackage{xspace}
\usepackage{cancel}
\usepackage{mathtools}
\usepackage{amsthm}
\makeatletter
\@ifundefined{thm}{\newtheorem{thm}{Thm}}{}
\@ifundefined{lema}{\newtheorem{lema}[thm]{Lemma}}{}
\makeatother
\title[Classical solutions to the Boltzmann equations for gas mixtu]{Classical solutions to the Boltzmann equations for gas mixture with unequal molecular masses}
\author{Gaofeng Wang, Weike Wang, Tianfang Wu}
\date{Source: arXiv:2601.21213v1 (CC BY 4.0)}
\begin{document}
\maketitle

\noindent\emph{Source and licence.} Classical solutions to the Boltzmann equations for gas mixture with unequal molecular masses. Version arXiv:2601.21213v1 (CC BY 4.0), distributed under the Creative Commons Attribution 4.0 International licence, as recorded in the arXiv licence metadata for this exact version and shown on the abstract page https://arxiv.org/abs/2601.21213v1. Full rights notice: licenses/arxiv-2601.21213-CC-BY-4.0.txt.\par\smallskip

\noindent\textit{Editorial scope.} Counted unit: the Lema whose source label is \texttt{222maxdel} in the article (source0.tex lines 1578--1586, source label \texttt{222maxdel}), delivered together with the article's own complete original proof of it (source0.tex lines 1587--1712, one \texttt{proof} environment of 126 lines). No mathematics is written, repaired or re-derived: statement and proof are the article's own text line for line. Required context retained uncounted: lem22k2ssig (lines 739--749); lem22cclssig (lines 823--825); CE2easyPT (lines 854--861); CE2hadPT (lines 854--861); caseA1wgams (lines 983--985); lem2230cgvb (lines 1087--1092); sec2CaseB1 (lines 1164--1166); sec2CaseB4 (lines 1226--1228). Every definition, display and label of those lines is retained unchanged. Omitted from this package and not redistributed: 1--738, 750--822, 826--853, 862--982, 986--1086, 1093--1163, 1167--1225, 1229--1577, 1713--2985. Nothing inside a retained excerpt is omitted or altered: every retained body is delivered exactly as the article's lines are written, so no omission receipt of any kind accompanies this package. Citations: the retained excerpts carry no citation command at all, so no bibliography entry of the article is transcribed and no reference is redistributed. Reference adaptation: none was needed and none was made; every reference inside the delivered unit resolves inside it, so no unresolved reference or citation is produced. Printed slips kept literal: no printed word, symbol, hypothesis or number of the counted statement or of its proof was altered. Layout and class: the article's own class is \texttt{amsart} with packages including hyperref, amssymb, mathrsfs, esint, graphicx, wrapfig, float, inputenc, amsmath, bm; the wrapper is the lane's pinned \texttt{amsart} editorial wrapper at 10pt on A4, which restates the theorem-like environments the retained text needs and adds the article's own macro definitions that the retained text needs (none), with extra wrapper packages (bm, bbm, dsfont, xspace, cancel, mathtools). The delivered numbering is the unit's own. Assets: no retained span contains a figure environment, a graphics-inclusion command, a PDF-page inclusion, a table or a TikZ picture; no image file is copied into this package, the package declares no asset dependency and no image file is redistributed with it. Licence: version arXiv:2601.21213v1 of this article is distributed under the Creative Commons Attribution 4.0 International licence, as recorded in the arXiv licence metadata for this exact version and shown on the abstract page https://arxiv.org/abs/2601.21213v1. A full rights notice is delivered at licenses/arxiv-2601.21213-CC-BY-4.0.txt. Characters: every retained excerpt is pure ASCII, so the delivered engine, XeLaTeX with its default Latin Modern text font, reports no missing character.
		\begin{equation}\label{lem22k2ssig}
			\begin{split}
				C\sum_{\beta=A,B} &\left(\iiint_{(\mathbb{R}^3)^2 \times \mathbb{S}^2}
				\sqrt{\mu^\beta(v_*)}\sqrt{\mu^{\beta}(v_*')}|v - v_*|^\gamma w^{2l} (1 - \chi)[f^\alpha_1(v')]^2
				d\omega dv_*dv\right)^{1/2}\\
				&\hspace{0.5cm}\times
				\left(\iiint_{(\mathbb{R}^3)^2 \times \mathbb{S}^2}
				\sqrt{\mu^\beta(v_*)}\sqrt{\mu^{\beta}(v_*')}|v - v_*|^\gamma w^{2l} (1 - \chi)[f^\alpha_2(v)]^2
				d\omega dv_*dv\right)^{1/2}
			\end{split}	
		\end{equation}	

		\begin{equation}\label{lem22cclssig}
		\left \langle  w^{2l}\mathbf{K}_{2,s}^{1-\chi}\mathbf{f}_1,\mathbf{f}_2\right \rangle_{L^{2}_{v}}=	\sum_{\alpha=A,B}\left \langle  w^{2l}K_{2,s}^{\alpha,1-\chi}\mathbf{f}_1,f^{\alpha}_2\right \rangle_{L^{2}_{v}} \leq C \varepsilon^{3+\gamma} \big|w^{l} \mathbf{f}_1\big|_\nu \big|w^{l} \mathbf{f}_2\big|_\nu
		\end{equation}

		\begin{align}
			\sum_{\beta=A,B}&\Bigg(\int_{|v|+|v_*|>m}\int_{\mathbb{S}^2} |u|^{\gamma} [\mu^{\beta}(u+v)]^{\frac{1}{2}}\chi(|u|)\big[\mu^{\beta}(v + u_{\perp}+\frac{m^{\beta}-m^{\alpha}}{m^{\alpha}+m^{\beta}}u_{\parallel})\big]^{\frac{1}{2}}\notag\\
			&\hspace{5.5cm}\times f^\alpha_1(v + \frac{2m^{\beta}}{m^{\alpha}+m^{\beta}}u_{\parallel})b^{ \alpha \beta}(\theta) \, du \, d\omega\Bigg) \label{CE2easyPT}\\
			& +\int_{|v|+|v_*|>m}\int_{\mathbb{S}^2} |u|^{\gamma} [\mu^{\beta}(u+v)]^{\frac{1}{2}}\chi(|u|)\big[\mu^{\alpha}(v + \frac{2m^{\beta}}{m^{\alpha}+m^{\beta}}u_{\parallel})\big]^{\frac{1}{2}}
			\hspace{1cm}(\alpha \neq \beta)\notag\\
			&\hspace{5cm}\times f^\beta_1(v + u_{\perp}+\frac{m^{\beta}-m^{\alpha}}{m^{\alpha}+m^{\beta}}u_{\parallel})b^{ \alpha \beta}(\theta) \, du \, d\omega \label{CE2hadPT}\\
			&+\int_{|v|+|v_*|>m}\int_{\mathbb{S}^2} |u|^{\gamma} [\mu^{\alpha}(u+v)]^{\frac{1}{2}}\chi(|u|)\big[\mu^{\alpha}(v + u_{\parallel})\big]^{\frac{1}{2}} f^\alpha_1(v + u_{\perp})b^{ \alpha \alpha}(\theta) \, du \, d\omega \label{CE2hadDGP}
		\end{align}

		\begin{equation}\label{caseA1wgams}
		C e^{-\frac{m^\beta}{16}(\frac{m^\alpha}{m^\alpha+m^\beta})^2|u_{\parallel}|^2}	\leq w^{\gamma-1}(v).
		\end{equation}

		\begin{eqnarray}\label{lem2230cgvb}
			\begin{split}
				\xi_{\parallel} &= \frac{\sqrt{m^\beta}+\sqrt{m^\alpha}}{2}v_{\parallel} + (\frac{\sqrt{m^\beta}}{2}+\frac{\sqrt{m^{\alpha}}m^{\beta}}{m^{\alpha}+m^{\beta}})u_{\parallel} \\
				\xi_{\perp} &= \frac{\sqrt{m^\beta}+\sqrt{m^\alpha}}{2}v_{\perp} +\frac{\sqrt{m^\beta}}{2}u_{\perp},
			\end{split}
		\end{eqnarray}

		\begin{equation}\label{sec2CaseB1}
			e^{-\frac{\bar{c}}{4}|v_{\parallel}|^2} \leq \frac{C}{m}, \hspace{1.9cm} e^{-\frac{\bar{c}}{4}[|v_{\parallel}|^2+|v_{\perp}|^2]} \leq C w^{\gamma}(v).
		\end{equation}

		\begin{equation}\label{sec2CaseB4}
			e^{-\frac{\bar{c}}{4}|u_{\perp}|^2} \leq \frac{C}{m}, \hspace{1.9cm} e^{-\frac{\bar{c}}{4}(|v_{\parallel}|^2+|v_{\perp}|^2)} \leq C w^{\gamma}(v).
		\end{equation}

	\begin{lema}\label{222maxdel}
		Let $\tilde{\beta}$ be  a multi-index and $l\geq0$. For any small positive $\eta$, there exists a  constant $C_{\eta}>0$ such that
		\begin{align}
			\big\langle w^{2l}  \partial_{\tilde{\beta}}[(\bm{\nu}:\mathbf{f})], \partial_{\tilde{\beta}} \mathbf{f} \big\rangle_{L_v^2} 
			&\geq  \big|w^{l}\partial_{\tilde{\beta}}\mathbf{f}\big|^2_{\nu}-\eta \sum_{|\tilde{\beta_1}|\leq|\tilde{\beta}|}\big|w^{l}\partial_{\tilde{\beta}_1}\mathbf{f}\big|^2_{\nu}-C_{\eta}\big|w^{l}\mathbf{f}\big|^2_{\nu};\label{lem24252tep}\\
			\big| \big\langle w^{2l} \partial_{\tilde{\beta}}[\mathbf{K} \mathbf{f}_1], \partial_{\tilde{\beta}} \mathbf{f}_2 \big\rangle\big| 
			&\leq \Bigg(\eta\sum_{|\tilde{\beta_1}|\leq|\tilde{\beta}|}|w^{l}\partial_{\tilde{\beta}_1}\mathbf{f}_1|_{\nu}+C_{\eta}|w^{l}\mathbf{f}_1|_{\nu}\Bigg)\big|w^{l}\partial_{\tilde{\beta}}\mathbf{f}_2\big|_{\nu}.\label{lem24253tep}
		\end{align}
	\end{lema}

	\begin{proof}
		By the chain rule of differentiation, 
		\begin{equation}\label{zzzdbbl}
			\big\langle w^{2l}  \partial_{\tilde{\beta}}[(\bm{\nu}:\mathbf{f})], \partial_{\tilde{\beta}} \mathbf{f} \big\rangle_{L_v^2}=\big\langle w^{2l}  [(\bm{\nu}:\partial_{\tilde{\beta}}\mathbf{f})], \partial_{\tilde{\beta}} \mathbf{f} \big\rangle_{L_v^2}+C_{\tilde{\beta}}^{\tilde{\beta_1}}\big\langle w^{2l}  [(\partial_{\tilde{\beta_1}}\bm{\nu}:\partial_{\tilde{\beta}-\tilde{\beta_1}}\mathbf{f})], \partial_{\tilde{\beta}} \mathbf{f} \big\rangle_{L_v^2}
		\end{equation}
		where $\tilde{\beta}\geq\tilde{\beta_1}$. Moreover, a direct calculation shows that
		\begin{equation}\label{nus24derivty}
			|\partial_{\tilde{\beta_1}}\bm{\nu}| \leq C (1+|v|)^{\gamma-1}.
		\end{equation}
		Then we obtain
		\begin{align*}
			\big\langle w^{2l}  [(\partial_{\tilde{\beta_1}}\bm{\nu}:\partial_{\tilde{\beta}-\tilde{\beta_1}}\mathbf{f})], \partial_{\tilde{\beta}} \mathbf{f} \big\rangle_{L_v^2} =& \big\langle w^{2l}  [(\partial_{\tilde{\beta_1}}\bm{\nu}:\partial_{\tilde{\beta}-\tilde{\beta_1}}\mathbf{f})], \mathbf{1}_{\{|v|>m\}}\partial_{\tilde{\beta}} \mathbf{f} \big\rangle_{L_v^2}\\
			&\hspace{0.4cm}+\big\langle w^{2l}  [(\partial_{\tilde{\beta_1}}\bm{\nu}:\partial_{\tilde{\beta}-\tilde{\beta_1}}\mathbf{f})], \mathbf{1}_{\{|v|\leq m\}}\partial_{\tilde{\beta}} \mathbf{f} \big\rangle_{L_v^2}.
		\end{align*}
		By \eqref{nus24derivty}, the part of $|v|\geq m$ is bounded by
		\begin{align*}
			\big\langle w^{2l}  [(\partial_{\tilde{\beta_1}}\bm{\nu}:\partial_{\tilde{\beta}-\tilde{\beta_1}}\mathbf{f})], \mathbf{1}_{\{|v|>m\}}\partial_{\tilde{\beta}} \mathbf{f} \big\rangle_{L_v^2} &\leq \frac{C}{m} \big|w^{l}\partial_{\tilde{\beta}-\tilde{\beta_1}}\mathbf{f}\big|_{\nu} \big|w^{l}\partial_{\tilde{\beta}}\mathbf{f}\big|_{\nu} \\
			& \leq  \frac{\eta}{2} \big|w^{l}\partial_{\tilde{\beta}-\tilde{\beta_1}}\mathbf{f}\big|_{\nu} \big|w^{l}\partial_{\tilde{\beta}}\mathbf{f}\big|_{\nu}.
		\end{align*}
		For the $m$ chosen above, by Cauchy-Schwarz inequality as well as the compact interpolation in the Sobolev space, the part of $|v|\leq m$ is bounded by
		\begin{align*}
			\big\langle w^{2l}  [(\partial_{\tilde{\beta_1}}\bm{\nu}:\partial_{\tilde{\beta}-\tilde{\beta_1}}\mathbf{f})], \mathbf{1}_{\{|v|\leq m\}}\partial_{\tilde{\beta}} \mathbf{f} \big\rangle_{L_v^2} \leq \frac{\eta}{2} \sum_{|\tilde{\beta_1}|=|\tilde{\beta}|} \big|w^{l}\partial_{\tilde{\beta_1}}\mathbf{f}\big|_{\nu}^2 + C_{\eta} \big|w^{l}\mathbf{f}\big|_{\nu}^2.
		\end{align*}
		Therefore, this concludes \eqref{lem24252tep}.
		
		Next we turn to the estimate for \eqref{lem24253tep}. From the decomposition $\mathbf{K}=\mathbf{K}_1+\mathbf{K}_2$, first we consider $\big\langle w^{2l} \partial_{\tilde{\beta}}\mathbf{K}_1 \mathbf{f}_1, \partial_{\tilde{\beta}}\mathbf{f}_2 \big\rangle_{L^{2}_{v}}$, which is bounded by
			\begin{align}
				 & \sum_{\alpha=A,B}\sum_{\beta=A,B} C_{\tilde{\beta}}^{\tilde{\beta_1}} \iiint_{\mathbb{S}^2\times(\mathbb{R}^3)^2} w^{2l}(v)|v_*|^\gamma  \partial_{\tilde{\beta_1}}\big(\sqrt{\mu^\alpha(v)} \sqrt{\mu^\beta(v_*+v)}\big)\notag\\
				 &\hspace{6cm}\times\partial_{\tilde{\beta}-\tilde{\beta_1}}f^\beta_1(v_*+v)\partial_{\tilde{\beta}}f^\alpha_2(v)b^{\alpha \beta}(\theta)\,d\omega \, dv_*\,dv \notag\\
				&\hspace{0.4cm}\leq\sum_{\alpha=A,B}\sum_{\beta=A,B} C_{\tilde{\beta}}^{\tilde{\beta_1}} \Bigg(\iiint_{\mathbb{S}^2\times(\mathbb{R}^3)^2} w^{2l}(v)|v_*|^\gamma  \partial_{\tilde{\beta_1}}\big(\sqrt{\mu^\alpha(v)} \sqrt{\mu^\beta(v_*+v)}\big)\label{lem24k11}\\
				&\hspace{5.5cm}\times\chi(|v_{*}|)\partial_{\tilde{\beta}-\tilde{\beta_1}}f^\beta_1(v_*+v)\partial_{\tilde{\beta}}f^\alpha_2(v)b^{\alpha \beta}(\theta)\,d\omega \, dv_*\,dv \notag\\
				&\hspace{1cm}+ \iiint_{\mathbb{S}^2\times(\mathbb{R}^3)^2} w^{2l}(v)|v_*|^\gamma  \partial_{\tilde{\beta_1}}\big(\sqrt{\mu^\alpha(v)} \sqrt{\mu^\beta(v_*+v)}\big)\label{lem24k12}\\
				&\hspace{5.2cm}\times\big[1-\chi(|v_{*}|)\big]\partial_{\tilde{\beta}-\tilde{\beta_1}}f^\beta_1(v_*+v)\partial_{\tilde{\beta}}f^\alpha_2(v)b^{\alpha \beta}(\theta)\,d\omega \, dv_*\,dv \Bigg)\notag,
			\end{align}
		where we have changed variable $v_*-v\rightarrow v_*$. Note that $\cos\theta=\frac{\omega\cdot(v-v_*)}{|v-v_*|}$ and $|b^{\alpha \beta}(\theta)|\leq C$, it yields that
		\begin{equation*}
			\Big|\partial_{\tilde{\beta_1}}\big(\sqrt{\mu^\alpha(v)} \sqrt{\mu^\beta(v_*+v)}\big)\Big|\leq C e^{-\frac{m_s}{8}(|v+v_*|^2+|v|^2)},
		\end{equation*}
		where we denote $m_s=\min\{m^A,m^B\}$.
		Applying a further change of variable and integrating $\omega$ over $\mathbb{S}^2$. \eqref{lem24k12} is bounded by
		\begin{align*}
			&C\sum_{\alpha=A,B}\sum_{\beta=A,B}\sum_{\tilde{\beta_1}}\iint_{(\mathbb{R}^3)^2} e^{-\frac{m_s}{8}(|v_*|^2+|v|^2)}|v_*-v|^\gamma w^{2l}(v) \\
			&\hspace{6cm}\times\big[1-\chi(|v-v_{*}|)\big]\big|\partial_{\tilde{\beta}-\tilde{\beta_1}}f^\beta_1(v_*)\partial_{\tilde{\beta}}f^\alpha_2(v)\big| \, dv_*\,dv \\
			&\hspace{0.5cm}\leq C\sum_{\alpha=A,B}\sum_{\beta=A,B}\sum_{\tilde{\beta_1}}\Bigg(\int_{\mathbb{R}^3} \int_{\mathbb{R}^3}e^{-\frac{m_s|v|^2}{8}}|v_*-v|^\gamma w^{2l}(v)\big[1-\chi(|v-v_{*}|)\big] dv \\
			&\hspace{8.2cm}\times e^{-\frac{m_s|v_*|^2}{8}}|\partial_{\tilde{\beta}-\tilde{\beta_1}}f^\beta_1(v_*)|^2 \, dv_*\,\Bigg)^{1/2}\\
			&\hspace{1cm} \times \Bigg(\int_{\mathbb{R}^3} \int_{\mathbb{R}^3}e^{-\frac{m_s|v_*|^2}{8}}|v_*-v|^\gamma \big[1-\chi(|v-v_{*}|)\big] dv_* \,e^{-\frac{m_s|v|^2}{8}}w^{2l}(v)|\partial_{\tilde{\beta}}f^\alpha_2(v)|^2 dv\Bigg)^{1/2}\\
			&\hspace{0.5cm}\leq  C \varepsilon^{3+\gamma}\sum_{\tilde{\beta_1}}\big|w^{l}\partial_{\tilde{\beta}-\tilde{\beta}_1}\mathbf{f}_1\big|_{\nu}\,\big|w^{l}\partial_{\tilde{\beta}}\mathbf{f}_2\big|_{\nu}\\ &\hspace{0.5cm}\leq \frac{\eta}{2}\Big(\sum_{\tilde{\beta}_1\leq\tilde{\beta}}\big|w^{l}\partial_{\tilde{\beta}_1}\mathbf{f}_1\big|_{\nu}\Big)\big|w^{l}\partial_{\tilde{\beta}}\mathbf{f}_2\big|_{\nu}
		\end{align*}
		For the small $\varepsilon>0$ chosen above, since $\chi_{\varepsilon}$ vanishes near the origin, an integration over $v_*$ in \eqref{lem24k11} yields
		\begin{align*}
			& \sum_{\alpha=A,B}\sum_{\beta=A,B} C_{\tilde{\beta}}^{\tilde{\beta_1}} \iiint_{\mathbb{S}^2\times(\mathbb{R}^3)^2} \partial_{\tilde{\beta}-\tilde{\beta_1}}^{v_*}\Big[|v_*|^\gamma \chi(|v_{*}|) \partial_{\tilde{\beta_1}}\big(\sqrt{\mu^\alpha(v)} \sqrt{\mu^\beta(v_*+v)}\big)\Big]\\
			&\hspace{6.5cm}\times w^{2l}(v)f^\beta_1(v_*+v)\partial_{\tilde{\beta}}f^\alpha_2(v)b^{\alpha \beta}(\theta)\,d\omega \, dv_*\,dv \\
			&\hspace{0.5cm}\leq \sum_{\alpha=A,B}\sum_{\beta=A,B}C_{\varepsilon} \iint_{(\mathbb{R}^3)^2} e^{-\frac{m_s}{8}(|v+v_*|^2+|v|^2)}|f^\beta_1(v_*+v)\partial_{\tilde{\beta}}f^\alpha_2(v)| dv_* \, dv \\
			&\hspace{0.5cm}\leq  \, C_{\varepsilon}\, \big|w^{l}\mathbf{f}_1\big|_{\nu} \big|w^{l}\partial_{\tilde{\beta}}\mathbf{f}_2\big|_{\nu}
		\end{align*}
		where we have used notation $\partial_{\tilde{\beta}-\tilde{\beta_1}}^{v_*}$ to denote the derivative with respect to $v_*$, in order to distinguish it from notation $\partial_{\tilde{\beta}-\tilde{\beta_1}}$ which represents the derivative with respect to $v$. Therefore, the estimate for the part of $\mathbf{K}_1$ can be derived.
		
		Now we consider the part for $\mathbf{K}_2$. $\mathbf{K}_2\mathbf{f}_1$ takes the form
		\begin{align}
			\sum_{\beta=A,B}&\Bigg(\iint_{\mathbb{S}^2\times\mathbb{R}^3} |u|^{\gamma} [\mu^{\beta}(u+v)]^{\frac{1}{2}}\big[1-\chi(|u|)\big]\big[\mu^{\beta}(v + u_{\perp}+\frac{m^{\beta}-m^{\alpha}}{m^{\alpha}+m^{\beta}}u_{\parallel})\big]^{\frac{1}{2}}\label{K2f1essigf11}\\
			&\hspace{5.5cm}\times f^\alpha_1(v + \frac{2m^{\beta}}{m^{\alpha}+m^{\beta}}u_{\parallel})b^{ \alpha \beta}(\theta) \, du \, d\omega \notag\\
			&\hspace{0.3cm}+ \iint_{\mathbb{S}^2\times\mathbb{R}^3} |u|^{\gamma} [\mu^{\beta}(u+v)]^{\frac{1}{2}}\big[1-\chi(|u|)\big]\big[\mu^{\alpha}(v + \frac{2m^{\beta}}{m^{\alpha}+m^{\beta}}u_{\parallel})\big]^{\frac{1}{2}}\label{K2f1essigf12}\\
			&\hspace{5cm}\times f^\beta_1(v + u_{\perp}+\frac{m^{\beta}-m^{\alpha}}{m^{\alpha}+m^{\beta}}u_{\parallel})b^{ \alpha \beta}(\theta) \, du \, d\omega \notag\\
			&\hspace{0.3cm}+\iint_{\mathbb{S}^2\times\mathbb{R}^3} |u|^{\gamma} [\mu^{\beta}(u+v)]^{\frac{1}{2}}\chi(|u|)\big[\mu^{\beta}(v + u_{\perp}+\frac{m^{\beta}-m^{\alpha}}{m^{\alpha}+m^{\beta}}u_{\parallel})\big]^{\frac{1}{2}}\label{K2f1etssigf13}\\
			&\hspace{5.5cm}\times f^\alpha_1(v + \frac{2m^{\beta}}{m^{\alpha}+m^{\beta}}u_{\parallel})b^{ \alpha \beta}(\theta) \, du \, d\omega \notag\\
			&\hspace{0.3cm}+ \iint_{\mathbb{S}^2\times\mathbb{R}^3} |u|^{\gamma} [\mu^{\beta}(u+v)]^{\frac{1}{2}}\chi(|u|)\big[\mu^{\alpha}(v + \frac{2m^{\beta}}{m^{\alpha}+m^{\beta}}u_{\parallel})\big]^{\frac{1}{2}}\label{K2f1ressigf14}\\
			&\hspace{5.3cm}\times f^\beta_1(v + u_{\perp}+\frac{m^{\beta}-m^{\alpha}}{m^{\alpha}+m^{\beta}}u_{\parallel})b^{ \alpha \beta}(\theta) \, du \, d\omega \Bigg). \notag
		\end{align}
		To compute $\partial_{\tilde{\beta}} \mathbf{K}_2\mathbf{f}_1$, differentiate the second term above directly with respect to $\beta$ using the product rule. By a further change of variables, this approach bypasses the need to compute derivatives of the singular kernel $|u - v|^\gamma$ in \eqref{K2f1essigf11} and \eqref{K2f1essigf12}. 
		Consequently, the term $\langle w^{2l} \partial_{\tilde{\beta}} \mathbf{K}_2 \mathbf{f}_1, \partial_{\tilde{\beta}} \mathbf{f}_2 \rangle$ can be decomposed as follows:
			\begin{align}
			&\sum_{\alpha=A,B}\sum_{\beta=A,B}\sum_{\tilde{\beta}_{i}}C_{\tilde{\beta}}\Bigg(\iiint_{\mathbb{S}^2\times(\mathbb{R}^3)^2} w^{2l}(v) |v-v_*|^{\gamma} \partial_{\tilde{\beta}_1}[\mu^{\beta}(v_*)]^{\frac{1}{2}}\big[1-\chi(|v-v_*|)\big]\partial_{\tilde{\beta}_2}[\mu^{\beta}(v_*')]^{\frac{1}{2}}\notag\\
			&\hspace{8.2cm}\times \partial_{\tilde{\beta}_3}f^\alpha_1(v')\partial_{\tilde{\beta}}f^\alpha_2(v)b^{ \alpha \beta}(\theta) \, du \,dv\, d\omega \label{sma248}\\
			&\hspace{0.3cm}+ \iiint_{\mathbb{S}^2\times(\mathbb{R}^3)^2} w^{2l}(v) |v-v_*|^{\gamma} \partial_{\tilde{\beta}_1}[\mu^{\beta}(v_*)]^{\frac{1}{2}}\big[1-\chi(|v-v_*|)\big]\partial_{\tilde{\beta}_2}[\mu^{\alpha}(v')]^{\frac{1}{2}}\notag\\
			&\hspace{9cm}\times \partial_{\tilde{\beta}_3}f^\beta(v_*')\partial_{\tilde{\beta}}f^\alpha_2(v)b^{ \alpha \beta}(\theta) \, dv_* \, dv\, d\omega \Bigg)\label{sma249}\\
			&\hspace{0.3cm}+\sum_{\alpha=A,B}\sum_{\beta=A,B}\Bigg(\int_{\mathbb{R}^3} \partial_{\tilde{\beta}}\Big[\iint_{\mathbb{S}^2\times\mathbb{R}^3}|v-v_*|^{\gamma} [\mu^{\beta}(v_*)]^{\frac{1}{2}}\chi(|v-v_*|)[\mu^{\beta}(v_*')]^{\frac{1}{2}} f^\alpha(v')b^{ \alpha \beta}(\theta) \, du \, d\omega \Big]\notag\\
				&\hspace{10.5cm}\times w^{2l}(v) \partial_{\tilde{\beta}}f^\alpha_2(v) \, dv\label{bga250}\\
			&\hspace{0.3cm}+ \int_{\mathbb{R}^3} \partial_{\tilde{\beta}}\Big( \iint_{\mathbb{S}^2\times\mathbb{R}^3} |v-v_*|^{\gamma} [\mu^{\beta}(v_*)]^{\frac{1}{2}}\chi(|v-v_*|)[\mu^{\alpha}(v')]^{\frac{1}{2}}
			 f^\beta(v_*')b^{ \alpha \beta}(\theta) \, du \, d\omega \Big)
			 w^{2l}(v) \partial_{\tilde{\beta}}f^\alpha_2(v) \, dv\Bigg), \label{bga251}
		\end{align}
		where $\tilde{\beta}=\tilde{\beta}_1+\tilde{\beta}_2+\tilde{\beta}_3$. Since 
		\begin{equation*}
			\big|\partial_{\tilde{\beta}}[\mu^{\beta}(\cdot)]^{\frac{1}{2}}\big| \leq C e^{-\frac{m_s|\cdot|^2}{8}}.
		\end{equation*}
		Following a similar procedure as in \eqref{lem22k2ssig} to \eqref{lem22cclssig}, the terms \eqref{sma248} and \eqref{sma249} are bounded by
		\begin{equation*}
			\varepsilon^{3+\gamma}\sum_{\tilde{\beta_1}}|w^{l}\partial_{\tilde{\beta}-\tilde{\beta}_1}\mathbf{f}_1|_{\nu}\,|w^{l}\partial_{\tilde{\beta}}\mathbf{f}_2|_{\nu} \leq \frac{\eta}{2}\Big(\sum_{\tilde{\beta}_1\leq\tilde{\beta}}|w^{l}\partial_{\tilde{\beta}_1}\mathbf{f}_1|_{\nu}\Big)|w^{l}\partial_{\tilde{\beta}}\mathbf{f}_2|_{\nu}.
		\end{equation*}
		For the small $\varepsilon>0$ above, \eqref{bga250}, \eqref{bga251}can be written as
		\begin{align}
		\sum_{\alpha,\beta}\sum_{\tilde{\beta}_{i}}&C_{\tilde{\beta}}	\Bigg( \int_{\mathbb{R}^3} \frac{1}{|u_{\parallel}|} \partial_{\tilde{\beta}_1}\big(e^{-\frac{m^\beta}{2}|\zeta_{\parallel}|^2}\big) e^{-\frac{m^\beta}{2}(\frac{m^\alpha}{m^\alpha+m^\beta})^2|u_{\parallel}|^2}\partial_{\tilde{\beta}_2}f^\alpha(v + \frac{2m^{\beta}}{m^{\alpha}+m^{\beta}}u_{\parallel}) \label{CE2BDSefsD1}\\
			&\hspace{0.2cm}\times \int_{\mathbb{R}^2} \partial_{\tilde{\beta}_1}\big(e^{-\frac{m^\beta}{2}|u_{\perp} + \zeta_{\perp}|^2}\big) \left[|u_{\parallel}|^2 + |u_{\perp}|^2\right]^{\frac{\gamma - 1}{2}} \chi\big(\sqrt{|u_{\parallel}|^2 + |u_{\perp}|^2}\big) \frac{b^{ \alpha \beta}(\theta)}{|\cos \theta|} \, du_{\perp} du_{\parallel}\notag\\
			 +&\int_{\mathbb{R}^3} \frac{1}{|u_{\parallel}|} e^{-\frac{1}{4}[|\xi_{\parallel}|^2+|\eta_{\parallel}|^2]}\int_{\mathbb{R}^2} \frac{b^{ \alpha \beta}(\theta)}{|\cos \theta|}\chi\big(\sqrt{|u_{\parallel}|^2 + |u_{\perp}|^2}\big)  \label{CE2BDfSHHD2}\\
			&\hspace{1.3cm} \times  e^{-\frac{1}{4}(|\xi_{\perp}|^2+|\eta_{\perp}|^2) } \left[|u_{\parallel}|^2 + |u_{\perp}|^2\right]^{\frac{\gamma - 1}{2}}f^\beta(v + u_{\perp}+\frac{m^{\beta}-m^{\alpha}}{m^{\alpha}+m^{\beta}}u_{\parallel})  \, du_{\perp} du_{\parallel} \Bigg)\notag.
		\end{align}
		By the expression of  $\eta_{\parallel},\xi_{\perp},\eta_{\perp},\xi_{\parallel}$ and $\zeta_{\perp}, \zeta_{\parallel}$, we deduce that
		\begin{equation*}
			|\nabla_v \eta_{\parallel}|+|\nabla_v \eta_{\perp}|+|\nabla_v \xi_{\parallel}|+|\nabla_v \xi_{\perp}|+|\nabla_v \zeta_{\parallel}|+|\nabla_v \zeta_{\perp}| \leq C,
		\end{equation*}
		which further gives
		\begin{align*}
			\big|\partial_{\tilde{\beta}_i}\big(e^{-\frac{m^\beta}{2}|\zeta_{\parallel}|^2}\big)\big|&\leq C e^{-\frac{m^\beta}{4}|\zeta_{\parallel}|^2},\\
			\big|\partial_{\tilde{\beta}_i}\big(e^{-\frac{m^\beta}{2}|u_{\perp} + \zeta_{\perp}|^2}\big)\big|&\leq C e^{-\frac{m^\beta}{4}|u_{\perp} + \zeta_{\perp}|^2},\\
			\big|\partial_{\tilde{\beta}_1}\big(e^{-\frac{1}{4}(|\xi_{\parallel}|^2+|\eta_{\parallel}|^2)}\big)\big|&\leq C e^{-\frac{1}{8}(|\xi_{\parallel}|^2+|\eta_{\parallel}|^2)},\\
			\big|\partial_{\tilde{\beta}_1}\big(e^{-\frac{1}{4}(|\xi_{\perp}|^2+|\eta_{\perp}|^2) }\big)\big|&\leq C e^{-\frac{1}{8}(|\xi_{\perp}|^2+|\eta_{\perp}|^2) }.
		\end{align*}
		Consequently, apart from the factors in the exponents, equation \eqref{CE2BDSefsD1},\eqref{CE2BDfSHHD2} is almost identical to \eqref{CE2easyPT},\eqref{CE2hadPT}. By choosing $m>0$ large enough, the integration domain is further partitioned into $(|v|+|v_*| \leq m)$ and $(|v|+|v_*| \geq m)$. 
		
		Applying analogous estimation techniques for  $\mathbf{K}_{2,r}^{\chi}\mathbf{f}_1$ as demonstrated in  \eqref{caseA1wgams}--\eqref{lem2230cgvb} and \eqref{sec2CaseB1}--\eqref{sec2CaseB4},   the contribution from the region $(|v|+|v_*| \geq m)$ can be bounded by
		\begin{equation*}
			\frac{C_{\varepsilon}}{m}\Bigg(\sum_{\tilde{\beta_1}}\big|w^{l}\partial_{\tilde{\beta}-\tilde{\beta}_1}\mathbf{f}_1\big|_{\nu}\Bigg)\,\big|w^{l}\partial_{\tilde{\beta}}\mathbf{f}_2\big|_{\nu} \leq \frac{\eta}{2}\Bigg(\sum_{\tilde{\beta}_1\leq\tilde{\beta}}\big|w^{l}\partial_{\tilde{\beta}_1}\mathbf{f}_1\big|_{\nu}\Bigg)\big|w^{l}\partial_{\tilde{\beta}}\mathbf{f}_2\big|_{\nu}.
		\end{equation*}
		Finally, we consider the remaining part, the operator  $\mathbf{K}_{2,z}^{\chi}$
			\begin{align*}
			&\sum_{\alpha,\beta}\sum_{\tilde{\beta}_{i}}C_{\tilde{\beta}}	\Bigg( \int_{\mathbb{R}^3} \frac{1}{|u_{\parallel}|} \partial_{\tilde{\beta}_1}\big(e^{-\frac{m^\beta}{2}|\zeta_{\parallel}|^2}\big) e^{-\frac{m^\beta}{2}(\frac{m^\alpha}{m^\alpha+m^\beta})^2|u_{\parallel}|^2}\partial_{\tilde{\beta}_2}f^\alpha(v + \frac{2m^{\beta}}{m^{\alpha}+m^{\beta}}u_{\parallel}) \\
			&\hspace{0.5cm}\times \int_{\mathbb{R}^2} \partial_{\tilde{\beta}_1}\big(e^{-\frac{m^\beta}{2}|u_{\perp} + \zeta_{\perp}|^2}\big) \left[|u_{\parallel}|^2 + |u_{\perp}|^2\right]^{\frac{\gamma - 1}{2}} \chi\big(\sqrt{|u_{\parallel}|^2 + |u_{\perp}|^2}\big) \frac{b^{ \alpha \beta}(\theta)}{|\cos \theta|}\mathbf{1}_{|v|+|v_*|\leq m} \, du_{\perp} du_{\parallel}\\
			&\hspace{0.2cm}+\int_{\mathbb{R}^3} \frac{1}{|u_{\parallel}|} e^{-\frac{1}{4}(|\xi_{\parallel}|^2+|\eta_{\parallel}|^2)}\int_{\mathbb{R}^2} \frac{b^{ \alpha \beta}(\theta)}{|\cos \theta|}\chi\big(\sqrt{|u_{\parallel}|^2 + |u_{\perp}|^2}\big)  \\
			& \hspace{0.5cm}\times  e^{-\frac{1}{4}(|\xi_{\perp}|^2+|\eta_{\perp}|^2) } \left[|u_{\parallel}|^2 + |u_{\perp}|^2\right]^{\frac{\gamma - 1}{2}}f^\beta(v + u_{\perp}+\frac{m^{\beta}-m^{\alpha}}{m^{\alpha}+m^{\beta}}u_{\parallel})\mathbf{1}_{|v|+|v_*|\leq m}  \, du_{\perp} du_{\parallel} \Bigg)
		\end{align*}
		is compact since $\frac{1}{|u_{\parallel}|} \in L^2_{loc}(\mathbb{R}^3)$. Hence, there exists a $C_{\eta}>0$ such that $\left \langle  w^{2l}\mathbf{K}_{2,r}^{\chi}\mathbf{f}_1,\mathbf{f}_2 \right \rangle_{\nu}$ is bounded by
		\begin{align*}
			 \Bigg(\frac{\eta}{4} \sum_{|\tilde{\beta_1}|\leq|\tilde{\beta}|} \big|w^{l}\partial_{\tilde{\beta_1}}\mathbf{f}\big|_{\nu}^2 + C_{\eta} \big|w^{l}\mathbf{f}\big|_{\nu}^2\Bigg)\big|w^{l}\partial_{\tilde{\beta}}\mathbf{f}_2\big|_{\nu}.
		\end{align*}
		Based on the above analysis, we have completed the proof of this lemma.
	\end{proof}

\end{document}
